% TEMPLATE-MILC-v3.tex - plantilla maestra para todas las guias MILC v3
% Ver scripts/generadores/build-guia-milc.py para la lista completa de
% claves esperadas. El script valida que no queden placeholders sin
% reemplazar antes de escribir el archivo final.

\documentclass[11pt,letterpaper]{article}
\usepackage[margin=1.75cm]{geometry}
\usepackage[table]{xcolor}
\usepackage{fontspec}
\usepackage[spanish]{babel}
\usepackage{tikz}
\usepackage[most]{tcolorbox}
\usepackage{tabularx}
\usepackage{array}
\usepackage{fancyhdr}
\usepackage{titlesec}
\usepackage{ragged2e}
\usepackage{enumitem}
\usepackage{microtype}

\setmainfont{Helvetica Neue}
\setsansfont{Helvetica Neue}
\setlength{\parindent}{0pt}
\setlength{\parskip}{6pt}
\hyphenpenalty=200
\exhyphenpenalty=200
\pretolerance=400
\tolerance=2000
\setlength{\emergencystretch}{1em}
\renewcommand{\arraystretch}{1.25}
\setlist{leftmargin=5mm,itemsep=1mm,topsep=1mm}
\newcolumntype{Y}{>{\RaggedRight\arraybackslash}X}

\definecolor{milcMagenta}{HTML}{E6007E}
\definecolor{milcVino}{HTML}{5A0038}
\definecolor{milcTurquesa}{HTML}{0093A5}
\definecolor{milcVerde}{HTML}{58A923}
\definecolor{milcMostaza}{HTML}{E5B400}
\definecolor{milcOcre}{HTML}{B86600}
\definecolor{milcNegro}{HTML}{241816}
\definecolor{milcGris}{HTML}{F7F7F4}
\definecolor{milcLinea}{HTML}{E8E2DD}
\definecolor{milcGradePrimary}{HTML}{84CC16}
\definecolor{milcGradeDark}{HTML}{4D7C0F}
\definecolor{milcGradeSoft}{HTML}{ECFCCB}
\definecolor{milcGradeAccent}{HTML}{CFE7FF}
\definecolor{milcGradeText}{HTML}{FFFFFF}
\color{milcNegro}

\pagestyle{fancy}
\fancyhf{}
\lhead{\small Tecnología e Informática · Grado 7}
\rhead{\small Guía 5 · MILC v3}
\cfoot{\small\thepage}
\renewcommand{\headrulewidth}{0pt}

\newtcolorbox{titlebox}[1]{
  enhanced,colback=#1,colframe=#1,arc=7mm,boxrule=0pt,
  left=7mm,right=7mm,top=3mm,bottom=3mm,
  before skip=8pt,after skip=8pt,
  fontupper=\sffamily\bfseries\Large\color{white}
}

\newtcolorbox{softbox}[3]{
  enhanced,colback=#2,colframe=#1,borderline west={4pt}{0pt}{#1},
  arc=2mm,boxrule=.4pt,left=4mm,right=4mm,top=3mm,bottom=3mm,
  before skip=6pt,after skip=8pt,title={#3},coltitle=white,
  colbacktitle=#1,fonttitle=\sffamily\bfseries,fontupper=\RaggedRight,
  attach boxed title to top left={xshift=3mm,yshift=-2mm},
  boxed title style={arc=2mm, boxrule=0pt}
}

\newtcolorbox{drawbox}[2]{
  enhanced,colback=white,colframe=#1,arc=4mm,boxrule=1pt,
  left=5mm,right=5mm,top=4mm,bottom=4mm,before skip=8pt,after skip=8pt,
  title={#2},coltitle=white,colbacktitle=#1,fonttitle=\sffamily\bfseries,
  fontupper=\RaggedRight,
  attach boxed title to top left={xshift=4mm,yshift=-2mm},
  boxed title style={arc=3mm, boxrule=0pt}
}

\newtcolorbox{infoband}[1]{
  enhanced,colback=#1!8,colframe=#1!50,arc=2mm,boxrule=.5pt,
  left=4mm,right=4mm,top=2mm,bottom=2mm,before skip=4pt,after skip=4pt,
  fontupper=\small\RaggedRight
}

% Puente narrativo entre secciones (contrato editorial Conéctate).
% Da continuidad: "Acabas de X. Ahora vas a Y porque Z."
% Visualmente: banda izquierda en color del grado, texto en itálica pequeña.
\newtcolorbox{puentebox}{
  enhanced,colback=milcGradeSoft,colframe=milcGradeDark,
  borderline west={3pt}{0pt}{milcGradeDark},
  arc=0mm,boxrule=0pt,left=5mm,right=4mm,top=2.5mm,bottom=2.5mm,
  before skip=4pt,after skip=8pt,
  fontupper=\itshape\small\color{milcNegro}\RaggedRight
}
\newcommand{\puente}[1]{%
  \begin{puentebox}%
    \textcolor{milcGradeDark}{\textbf{\textrightarrow}}\hspace{2mm}#1%
  \end{puentebox}%
}

\newcommand{\linead}[1]{\noindent\color{milcLinea}\rule{#1}{0.5pt}\color{milcNegro}}
\newcommand{\respuesta}[1]{\par\vspace{2mm}\foreach \n in {1,...,#1}{\noindent\color{milcLinea}\rule{\linewidth}{0.5pt}\par\vspace{5mm}}\color{milcNegro}}
\newcommand{\checkbox}{\textcolor{milcGradeDark}{\fbox{\phantom{x}}}\hspace{5pt}}

\newcommand{\phasecard}[4]{
\begin{tcolorbox}[enhanced,colback=#3,colframe=#2,arc=3mm,boxrule=.5pt,height=3.05cm,valign=center,halign=center,left=1.5mm,right=1.5mm,top=2mm,bottom=2mm]
{\sffamily\bfseries\fontsize{22}{24}\selectfont\textcolor{#2}{#1}}\par
{\sffamily\bfseries\small #4}
\end{tcolorbox}
}

\begin{document}
\thispagestyle{empty}
\begin{tikzpicture}[remember picture,overlay]
  \fill[white] (current page.south west) rectangle (current page.north east);
  \path[fill=milcGradePrimary,rounded corners=16mm]
    ([xshift=.55cm,yshift=-.55cm]current page.north west) rectangle
    ([xshift=-.55cm,yshift=3.65cm]current page.south east);

  \node[anchor=north west,text=milcGradeText] at ([xshift=2.0cm,yshift=-1.85cm]current page.north west)
    {\sffamily\bfseries\fontsize{14}{16}\selectfont GRADO};
  \node[anchor=north west,text=milcGradeText,opacity=.82] at ([xshift=2.0cm,yshift=-2.75cm]current page.north west)
    {\sffamily\bfseries\fontsize{9}{11}\selectfont SERIE GUÍAS · TECNOLOGÍA E INFORMÁTICA};

  \begin{scope}[shift={([xshift=-3.05cm,yshift=-2.55cm]current page.north east)}]
    \fill[white,opacity=.80] (-.62,-.52) rectangle (.62,.52);
    \draw[milcGradeText,line width=1pt] (-.62,-.52) rectangle (.62,.52);
    \fill[milcGradeDark] (-.42,-.38) rectangle (-.22,.06);
    \fill[milcGradeAccent] (-.10,-.38) rectangle (.10,.24);
    \fill[milcGradePrimary!60!black] (.22,-.38) rectangle (.42,.38);
  \end{scope}

  \node[anchor=west,text=milcGradeText] at ([xshift=1.9cm,yshift=-12.85cm]current page.north west)
    {\sffamily\bfseries\fontsize{124}{124}\selectfont 7};
  \draw[milcGradeText,line width=1.7pt]
    ([xshift=4.52cm,yshift=-11.68cm]current page.north west) circle (.13cm);

  \node[anchor=west,text=milcGradeText,text width=8.7cm,align=left] at ([xshift=10.35cm,yshift=-11.6cm]current page.north west)
    {
     {\sffamily\bfseries\fontsize{19}{22}\selectfont Compartir y\\permisos ---\\abrir mi\\taller a\\invitados}\\[4mm]
     {\sffamily\bfseries\fontsize{23}{26}\selectfont Guía 5-7-TIC}\\[2mm]
     {\sffamily\fontsize{13}{16}\selectfont Período 1 · Trabajo colaborativo en la nube}\\[5mm]
     {\sffamily\bfseries\fontsize{11}{13}\selectfont Institución Educativa Sor María Juliana}\\[1mm]
     {\sffamily\fontsize{10}{12}\selectfont Autor: Dr. Álvaro Cárdenas Orozco}
    };

  \path[fill=milcGradeSoft,rounded corners=7mm]
    ([xshift=1.35cm,yshift=.85cm]current page.south west) rectangle
    ([xshift=-1.35cm,yshift=4.25cm]current page.south east);
  \draw[milcGradeDark,line width=.65pt,rounded corners=7mm]
    ([xshift=1.35cm,yshift=.85cm]current page.south west) rectangle
    ([xshift=-1.35cm,yshift=4.25cm]current page.south east);
  \node[anchor=center,text=milcNegro,text width=17.0cm,align=center] at ([yshift=2.55cm]current page.south)
    {\sffamily\fontsize{10}{12}\selectfont Metodología MILC · Apertura ancestral · Pensamiento computacional · Triángulo Dussel-Estoicismo-Floridi\\
    \textcolor{milcGradeDark}{\bfseries Producto:} Uno de tus archivos de la S4 (Word, Excel o PowerPoint) \textbf{compartido con 2 compañeros} de la clase, con permisos \textbf{distintos a cada uno}: uno \emph{solo puede ver}, el otro \emph{puede editar}. Más \textbf{en el cuaderno}: tabla de los 3 niveles de permisos (puede ver, comentar, editar) con ejemplos de cuándo usar cada uno + las 5 reglas del compartir responsable + captura impresa o boceto de la pantalla \emph{``Compartir''}.};
\end{tikzpicture}
\mbox{}
\newpage

\section*{Apertura: Compartir y permisos --- abrir mi taller a invitados con criterio}

\begin{softbox}{milcGradeDark}{milcGradeSoft}{Saber ancestral}
\textbf{Cuando un maestro artesano de Cartago abría su taller a un aprendiz, decidía exactamente qué podía hacer y qué no.} En el taller de don Lucho el relojero, no todos los aprendices podían tocar las mismas cosas. \textbf{Al recién llegado}, don Lucho le decía: \emph{``Mire, no toque''}. Era nivel \textbf{lector}. \textbf{Al aprendiz de 6 meses}, le decía: \emph{``Mire, dígame si ve algo raro, pero todavía no arregle''}. Era nivel \textbf{comentarista}. \textbf{Al aprendiz de 2 años}, ya: \emph{``Tome este reloj y repárelo, después le reviso''}. Era nivel \textbf{editor}. \textbf{¿Por qué tantos niveles?} Porque \emph{abrir el taller no es regalar el taller}. Cada nivel de acceso protegía el oficio --- un aprendiz inexperto podía dañar un reloj caro tocando lo equivocado. Hoy en Microsoft 365 pasa exactamente igual: cuando compartes un archivo, escoges \textbf{el nivel de permisos}. Saber escogerlo es la diferencia entre colaborar bien y armar un desorden.
\end{softbox}

\begin{softbox}{milcTurquesa}{milcGris}{Saber contemporáneo}
\textbf{Compartir} un archivo en OneDrive significa darle acceso a otra persona desde su cuenta. NO se copia el archivo: el archivo queda donde está, pero ahora otra persona también lo puede abrir. \textbf{Hay 3 niveles de permisos:} \textbf{(1) Puede ver} (\emph{lector}): solo lee el archivo, no puede cambiar nada. Ideal para entregar trabajos terminados al profe. \textbf{(2) Puede comentar} (\emph{comentarista}): puede leer y agregar comentarios al margen, pero no cambia el texto. Ideal para revisión. \textbf{(3) Puede editar} (\emph{editor}): puede leer, comentar y cambiar el archivo libremente. Ideal para coautoría real. \textbf{Cómo se comparte:} botón \emph{``Compartir''} (esquina superior derecha) → escribir correo del invitado + elegir nivel → \emph{``Enviar''}. El otro recibe un correo con el enlace y abre el archivo. \emph{Tú decides qué puede hacer él}.
\end{softbox}

\begin{softbox}{milcMostaza}{milcGris}{Pregunta-puente}
\textit{Si compartes tu archivo de Word con tu mejor amigo dándole \emph{``puede editar''} y él, sin querer, borra una página, ¿\textbf{se pierde para siempre} la información? ¿Cómo deberías haber compartido en lugar de eso?}
\end{softbox}

\begin{softbox}{milcVerde}{milcGris}{Saber y saber hacer}
Al terminar la clase: (1) podrás \textbf{identificar} los 3 niveles de permisos; (2) sabrás \textbf{aplicar} el nivel correcto a 5 situaciones; (3) podrás \textbf{evaluar} las 5 reglas del compartir responsable; (4) habrás \textbf{compartido} 1 archivo con 2 compañeros con permisos distintos.
\end{softbox}

\small
\begin{tabularx}{\textwidth}{>{\bfseries}p{3.0cm}Y}
\rowcolor{milcGradePrimary!12}Autor & Dr. Álvaro Cárdenas Orozco\\
DBA & Utiliza herramientas digitales para trabajar de manera colaborativa con otros (MEN, T\&I 7°)\\
Referentes & Saber ancestral de la minga · Microsoft 365 y OneDrive · Coautoría asincrónica\\
Producto final & Uno de tus archivos de la S4 (Word, Excel o PowerPoint) \textbf{compartido con 2 compañeros} de la clase, con permisos \textbf{distintos a cada uno}: uno \emph{solo puede ver}, el otro \emph{puede editar}. Más \textbf{en el cuaderno}: tabla de los 3 niveles de permisos (puede ver, comentar, editar) con ejemplos de cuándo usar cada uno + las 5 reglas del compartir responsable + captura impresa o boceto de la pantalla \emph{``Compartir''}.\\
\end{tabularx}
\normalsize

\puente{Hoy haces 4 cosas: aprendes los 3 niveles de permisos, los aplicas a 5 situaciones, conoces las 5 reglas del compartir, compartes un archivo real.}

\begin{titlebox}{milcGradeDark}
Ruta MILC: así vamos a aprender
\end{titlebox}
\begin{tabularx}{\textwidth}{>{\centering\arraybackslash}X>{\centering\arraybackslash}X>{\centering\arraybackslash}X>{\centering\arraybackslash}X}
\phasecard{1}{milcVerde}{milcGris}{Escuta\\[-1mm]\small Observo, escucho y pregunto.} &
\phasecard{2}{milcTurquesa}{milcGris}{Sistematización\\[-1mm]\small Pensamiento computacional.} &
\phasecard{3}{milcMagenta}{milcGris}{Praxis\\[-1mm]\small Construyo y aplico.} &
\phasecard{4}{milcVino}{milcGris}{Evaluación\\[-1mm]\small Reflexión liberadora.}
\end{tabularx}


\puente{Empezamos con 5 situaciones para escoger el nivel correcto.}

\begin{titlebox}{milcVerde}
Fase 1 · Escuta
\end{titlebox}

\begin{softbox}{milcVerde}{milcGris}{Escena para escuchar}
\textbf{Actividad 1 · IDENTIFICA --- 5 situaciones, 3 niveles} (10 min · individual). Lee estas 5 situaciones y escribe en el cuaderno \textbf{qué nivel de permiso (V/C/E)} darías. \emph{V = Puede ver. C = Puede comentar. E = Puede editar}. \textbf{(1) Le mandas al profe la tarea final terminada para que la califique.} \textbf{(2) Le pides a un compañero que revise tu ensayo y te diga si hay errores antes de entregar.} \textbf{(3) Vas a hacer un trabajo en grupo con 3 compañeros y todos van a escribir partes distintas.} \textbf{(4) Le mandas a tu mamá la presentación para que la vea antes de la exposición.} \textbf{(5) Le compartes con un amigo cercano tu archivo Word para que lo termine en tu lugar mientras tú estudias mate.}
\end{softbox}

\checkbox Leí las 5 situaciones con calma.\\
\checkbox Asigné V, C o E a cada una.\\
\checkbox La situación 5 me hizo pensar (¿es buena idea?).

\begin{infoband}{milcVerde}
\textbf{Lo que va al cuaderno (Actividad 1 · IDENTIFICA).} Título: «Actividad 1 --- 5 situaciones, 3 niveles». \textbf{Formato:} lista numerada del 1 al 5 con V/C/E. \textbf{Extensión:} 5 líneas. \textbf{Sabes que terminaste cuando} tienes una decisión por cada situación y podrías justificarla.
\end{infoband}

\puente{Después aprendes los 3 niveles + 5 reglas del compartir responsable.}

\begin{titlebox}{milcTurquesa}
Fase 2 · Sistematización con pensamiento computacional
\end{titlebox}

\begin{softbox}{milcTurquesa}{milcGris}{Cuatro pilares aplicados al tema}
Tu intuición se afila ahora. Aprende los \textbf{3 niveles exactos} con ejemplos profesionales y las 5 reglas para no hacer mal uso del compartir.
\end{softbox}

\small
\begin{tabularx}{\textwidth}{>{\bfseries\RaggedRight\arraybackslash}p{3.6cm}Y}
\rowcolor{milcTurquesa!90}\color{white}Pilar & \color{white}\bfseries Aplicación al tema\\
1. Descomposición & \textbf{Nivel 1 --- Puede ver (\emph{Read-only}).} El invitado abre el archivo pero \textbf{no puede modificar nada}. Si intenta escribir, sale un mensaje \emph{``solo lectura''}. \textbf{Cuándo usarlo:} entregar trabajos terminados al profe, mostrar archivos a familiares, compartir documentación oficial. \textbf{Pista profesional:} es el nivel más seguro. Si dudas, empieza por aquí. Si después necesitan más permisos, lo cambias.\\
2. Reconocimiento de patrones & \textbf{Nivel 2 --- Puede comentar.} El invitado puede leer y agregar \textbf{comentarios al margen} (como pegatinas amarillas digitales), pero \textbf{NO puede cambiar el texto}. Sus comentarios aparecen en el lateral del documento; tú decides si los aceptas o no. \textbf{Cuándo usarlo:} \emph{revisión de trabajos por compañeros o profe}, \emph{retroalimentación antes de entregar}, \emph{discusión sobre un texto sin que se cambie}. Lo verás más a fondo en la S7.\\
3. Abstracción & \textbf{Nivel 3 --- Puede editar.} El invitado puede \textbf{cambiar el documento libremente}: escribir, borrar, mover, agregar imágenes. Es el nivel de la \emph{coautoría real}. \textbf{Cuándo usarlo:} trabajos en grupo donde varios escriben, proyectos colaborativos, bitácoras compartidas. \textbf{Cuidado importante:} con editor, el otro puede dañar tu trabajo (sin querer o queriendo). Por eso da editor solo a personas que conozcas bien y solo cuando realmente sea necesario.\\
4. Algoritmo conceptual & \textbf{Las 5 reglas del compartir responsable.} (1) \textbf{El menor permiso necesario:} si solo necesitas que vean, no des editor. Es como prestar el carro: das las llaves solo si es necesario. (2) \textbf{Verifica el correo antes de enviar:} un punto de más y el archivo va a un desconocido. (3) \textbf{Pon plazo de expiración cuando sea posible:} \emph{``acceso solo hasta el viernes''}. (4) \textbf{Quita acceso cuando ya no se necesite:} en el panel de Compartir puedes ver quién tiene acceso y revocarlo. (5) \textbf{No compartas archivos con datos personales sensibles:} contraseñas, números de cédula, fotos íntimas. Esos no van a compartirse.\\
\end{tabularx}
\normalsize

\begin{drawbox}{milcTurquesa}{3 niveles + 5 reglas}
\textbf{3 niveles:} \\[1mm]
\textbf{1.} Puede VER --- solo lectura. \\[1mm]
\textbf{2.} Puede COMENTAR --- lectura + notas al margen. \\[1mm]
\textbf{3.} Puede EDITAR --- cambio libre. \\[1mm]
\textbf{5 reglas:} menor permiso · verifica correo · pon expiración · quita acceso · no datos sensibles.
\end{drawbox}

\begin{softbox}{milcMostaza}{milcGris}{Errores comunes y su corrección}
\textbf{Errores frecuentes al compartir.} (1) \emph{Dar editor a todo el mundo por comodidad} --- alguien va a dañar el archivo. (2) \emph{Mandar el enlace por WhatsApp sin verificar el destinatario} --- llega a alguien equivocado. (3) \emph{Olvidar revocar acceso cuando un compañero ya no es del equipo} --- sigue viendo cosas privadas. (4) \emph{Compartir trabajos terminados como editor en lugar de viewer} --- el profe puede modificar tu tarea por error. (5) \emph{No verificar que el compañero RECIBIÓ y ABRIÓ el archivo} --- te enteras del problema cuando ya pasó el plazo.
\end{softbox}

\begin{infoband}{milcTurquesa}
\textbf{Lo que va al cuaderno (Actividad 2 · EXPLICA).} Título: «Actividad 2 --- 3 niveles + 5 reglas». \textbf{Formato:} bloque A con tabla de 3 filas, 2 columnas (nivel + cuándo usarlo + ejemplo). Bloque B con lista de 5 reglas. \textbf{Extensión:} 8 líneas. \textbf{Sabes que terminaste cuando} podrías recomendar el nivel correcto a un compañero que dude.
\end{infoband}

\puente{Con todo claro, abres OneDrive y compartes uno de tus archivos con 2 compañeros.}

\begin{titlebox}{milcMagenta}
Fase 3 · Praxis con pensamiento computacional
\end{titlebox}

\begin{softbox}{milcMagenta}{milcGris}{Construyendo el producto con los cuatro pilares}
\textbf{Actividad 3 · APLICA --- Compartir un archivo real con 2 compañeros} (28 min · individual + interacción con compañeros). Vas a tomar uno de tus archivos de la S4 (Word, Excel o PowerPoint --- el que más te gustó) y compartirlo con 2 compañeros de la clase con permisos distintos.
\end{softbox}

\small
\begin{tabularx}{\textwidth}{>{\bfseries\RaggedRight\arraybackslash}p{3.6cm}Y}
\rowcolor{milcMagenta!90}\color{white}Pilar & \color{white}\bfseries Aplicación al producto\\
Descomposición & \textbf{Paso 1 --- Escoge el archivo y los 2 compañeros.} En el cuaderno escribe: \emph{``Archivo que voy a compartir: \_\_\_''}. \emph{``Compañero 1 (va a ser viewer): \_\_\_''}. \emph{``Compañero 2 (va a ser editor): \_\_\_''}. Acuerda con ellos qué hacer (que abran el archivo y verifiquen sus permisos).\\
Patrones & \textbf{Paso 2 --- Abre el archivo y haz clic en Compartir.} Entra a OneDrive, abre tu archivo. En la esquina superior derecha hay un botón \textbf{``Compartir''} (con un ícono de personas). Haz clic.\\
Abstracción & \textbf{Paso 3 --- Configura el primer compartir (viewer).} Aparece una ventana. En el campo de \emph{``Para''} escribe el correo institucional de tu Compañero 1 (algo como \texttt{nombre1@sormariajuliana.edu.co}). Donde dice \emph{``Cualquier persona con el vínculo puede editar''} (o similar), haz clic y cambia a \textbf{``Puede ver''}. Escribe un mensaje corto: \emph{``Hola [nombre], te comparto este archivo para que lo veas. No tienes que editar nada, solo verlo.''}. Pulsa \emph{Enviar}.\\
Algoritmo de armado & \textbf{Paso 4 --- Configura el segundo compartir (editor).} Vuelve a Compartir. Escribe correo de Compañero 2. Cambia el nivel a \textbf{``Puede editar''}. Mensaje: \emph{``Hola [nombre], te comparto este para que agregues 1 línea al final. Es práctica de coautoría.''}. Enviar.\\
\end{tabularx}
\normalsize

\begin{drawbox}{milcMagenta}{Antes de cerrar la S5}
\checkbox Compartí 1 archivo con 2 compañeros. \\[1mm]
\checkbox El Compañero 1 tiene 'Puede ver' (no edita). \\[1mm]
\checkbox El Compañero 2 tiene 'Puede editar' (sí cambia el archivo). \\[1mm]
\checkbox Verifiqué con los compañeros que los permisos funcionan. \\[1mm]
\checkbox Boceto de la pantalla Compartir está en el cuaderno. \\[1mm]
\checkbox La reflexión personal está firmada.
\end{drawbox}

\begin{softbox}{milcMagenta}{milcGris}{Plantilla de trabajo}
\textbf{Estructura.} \textit{Paso 1:} elegir archivo + 2 compañeros. \textit{Paso 2:} botón Compartir (esquina superior derecha). \textit{Pasos 3 y 4:} compartir con 2 niveles distintos. \textit{Paso 5:} verificar con los compañeros que funcionan. \textit{Paso 6:} boceto + reflexión + firma.
\end{softbox}

\begin{infoband}{milcMagenta}
\textbf{Lo que va al cuaderno (Actividad 3 · APLICA).} Título: «Actividad 3 --- Mi primer compartir real». \textbf{Formato:} datos del compartir + boceto de pantalla + reflexión. \textbf{Extensión:} 1 página. \textbf{Sabes que terminaste cuando} tus 2 compañeros tienen acceso al archivo con sus permisos respectivos.
\end{infoband}

\puente{El producto es ese archivo compartido + tabla de permisos en cuaderno.}

\begin{titlebox}{milcMostaza}
Producto de la sesión
\end{titlebox}
\textbf{Archivo compartido con 2 compañeros · 2 niveles distintos · verificado}.
\begin{softbox}{milcMostaza}{milcGris}{Criterios de calidad}
(1) Hay 1 archivo compartido con 2 compañeros. (2) Los 2 compañeros tienen niveles DISTINTOS (uno viewer, uno editor). (3) Los compañeros confirmaron que pueden abrir con su nivel correcto. (4) Boceto de la pantalla Compartir está en el cuaderno. (5) La reflexión personal está firmada.
\end{softbox}

\puente{Antes de salir, verifica con los compañeros que ellos pueden abrir el archivo con el permiso correcto.}

\begin{titlebox}{milcVino}
Fase 4 · Evaluación liberadora — cinco dimensiones
\end{titlebox}

\begin{softbox}{milcVino}{milcGris}{Sentido de la evaluación}
Evaluar no es solo revisar una nota. Es reconocer qué aprendí, qué cambió en mí, qué emociones manejé, cómo me relaciono con la ciudad y la comunidad, y qué le dejaré a quien viene detrás.
\end{softbox}

\textbf{Mi reflexión liberadora en cinco dimensiones.} Completa cada frase iniciada:

\small
\begin{tabularx}{\textwidth}{>{\bfseries\RaggedRight\arraybackslash}p{3.4cm}Y>{\RaggedRight\arraybackslash}p{4.6cm}}
\rowcolor{milcVino}\color{white}Dimensión & \color{white}\bfseries Mi respuesta (frame starter) & \color{white}\bfseries Algo que descubrí\\
1. Personal & Lo que más cambió en mí fue \linead{6.5cm} & \linead{4.2cm}\\
2. Emocional & La emoción más fuerte fue \linead{2.5cm} y la regulé \linead{3cm} & \linead{4.2cm}\\
3. Ciudadana & En democracia esto importa porque \linead{6cm} & \linead{4.2cm}\\
4. Local (Cartago) & En mi barrio sería útil para \linead{6.5cm} & \linead{4.2cm}\\
5. Intergeneracional & A mi abuela/abuelo le diría \linead{6.5cm} & \linead{4.2cm}\\
\end{tabularx}
\normalsize

\begin{infoband}{milcVino}
\textbf{Para la próxima sustentación.} Vuelve a esta tabla en seis meses. Verás que las respuestas cambian — no porque lo aprendido cambie, sino porque tú cambias. La evaluación liberadora es una conversación contigo a través del tiempo.
\end{infoband}


\puente{Cerramos con 3 voces que te ayudan a entender por qué dar acceso con criterio es respeto al trabajo.}

\begin{titlebox}{milcGradeDark}
Triángulo de pensamiento · cierre filosófico
\end{titlebox}

\begin{softbox}{milcVino}{milcGris}{Dussel · lente del nosotros}
\textit{````Compartir bien es escoger con criterio a quién das qué. No es dar todo a todos: eso no es generosidad, es descuido.''''}

Hay quien cree que \emph{``ser generoso''} es dar a todos el máximo acceso. Falso. \textbf{Generosidad sin criterio es descuido}: das editor a alguien que no necesita editor, y termina dañándose el trabajo. \emph{El verdadero compartir generoso es escoger el nivel correcto para cada persona}: el viewer a quien debe ver, el editor a quien debe escribir. Don Lucho lo entendía con sus aprendices. Tú lo entiendes hoy con tus archivos.

\textbf{¿Comparto con criterio o por comodidad?}
\end{softbox}
\linead{\linewidth}\\[1mm]
\linead{\linewidth}

\begin{softbox}{milcMostaza}{milcGris}{Estoicismo · Marco Aurelio (emperador que decidía con cuidado) · cuidado interior}
\textit{````Antes de dar, pregunta: ¿qué necesita esta persona específicamente? No lo que tú te imaginas; lo que ella requiere para hacer su parte.''''}

Compartir bien empieza por una pregunta: \textbf{¿qué necesita exactamente la otra persona?}. Si solo necesita ver tu trabajo, dale viewer. Si necesita agregar comentarios, comentarista. Si va a trabajar contigo en serio, editor. \emph{Cada nivel para cada necesidad}. Es como pedir herramientas en un taller: el aprendiz no se lleva el martillo si solo necesita la lupa.

\textbf{¿Pienso en lo que la otra persona necesita o en lo que es más rápido para mí?}
\end{softbox}
\linead{\linewidth}\\[1mm]
\linead{\linewidth}

\begin{softbox}{milcTurquesa}{milcGris}{Floridi · lente de la infoesfera}
\textit{````La diferencia entre un equipo digital que funciona y uno que no es la calidad de los permisos. Acceso fino = equipo eficaz.''''}

Los equipos digitales que \emph{funcionan bien} tienen un patrón: \textbf{cada persona tiene exactamente el acceso que necesita} (ni más, ni menos). Los que \emph{funcionan mal} suelen tener todos editores en todo, o todos viewers en todo. \emph{Trabajar bien empieza por configurar bien los permisos al inicio del proyecto}. Hoy aprendiste ese hábito. Te servirá toda la vida adulta.

\textbf{Si formara un equipo de trabajo, ¿sabría asignar permisos finos a cada integrante?}
\end{softbox}
\linead{\linewidth}\\[1mm]
\linead{\linewidth}

\begin{titlebox}{milcVino}
Compromiso final
\end{titlebox}

\small
\begin{tabularx}{\textwidth}{>{\RaggedRight\arraybackslash}p{3.0cm}YYY}
\rowcolor{milcVino}\color{white}\bfseries Criterio & \color{white}\bfseries Lo logré & \color{white}\bfseries En proceso & \color{white}\bfseries Necesito apoyo\\
Comprensión del tema & Explico ideas centrales con mis palabras. & Reconozco ideas pero falta precisión. & Necesito repasar conceptos base.\\
Pensamiento computacional & Apliqué los 4 pilares al diseñar mi producto. & Apliqué algunos pilares. & Necesito releer la fase 2.\\
Aplicación y producto & Mi producto cumple los criterios y se puede revisar. & Producto funcional con ajustes pendientes. & Aún en construcción.\\
Reflexión 5 dimensiones & Respondí cada dimensión con honestidad. & Respondí algunas con detalle. & Mis respuestas son muy generales.\\
Triángulo de pensamiento & Conecté las 3 voces con mi trabajo real. & Conecté al menos una. & No relaciono las voces con mi proceso.\\
\end{tabularx}
\normalsize

\textbf{Hoy entendí que:}\\[1mm]
\linead{\linewidth}\\[1mm]
\linead{\linewidth}

\textbf{Mi compromiso para la próxima sesión es:}\\[1mm]
\linead{\linewidth}\\[1mm]
\linead{\linewidth}

\begin{softbox}{milcMostaza}{milcGris}{Cita de despedida · Marco Aurelio}
\textit{``El obstáculo en el camino se vuelve el camino. Lo que estorba al camino se vuelve camino.''}\\[1mm]
Cada error de hoy, bien observado, se convierte en pieza del camino que pisarás mañana.
\end{softbox}

\small
\begin{tabularx}{\textwidth}{>{\bfseries}p{3.6cm}Y}
\rowcolor{milcGradeDark!12}Nombre completo: & \linead{10cm}\\
Fecha: & \linead{4cm} \quad Curso/grupo: \linead{4cm}\\
Mi firma: & \linead{9cm}\\
\end{tabularx}
\normalsize

\begin{infoband}{milcGradeDark}
\textbf{Para la próxima sesión.} Esta semana, intenta compartir 1 archivo más con un familiar (papá, mamá, abuelos) usando \emph{``Puede ver''}: enséñales tu trabajo escolar. Pasa de \emph{``mostrar pantalla''} a \emph{``compartir vínculo''} --- es paso pequeño pero importante. La próxima clase aprendes \textbf{coautoría asincrónica}: cómo 3 personas escriben al mismo tiempo en el mismo documento sin chocar.
\end{infoband}

\end{document}
